%% LyX 2.4.4 created this file.  For more info, see https://www.lyx.org/.
%% Do not edit unless you really know what you are doing.
\documentclass[english]{article}
\usepackage[T1]{fontenc}
\usepackage[latin9]{inputenc}
\usepackage{enumitem}
\usepackage{amsmath}
\usepackage{amsthm}
\usepackage{amssymb}
\usepackage{geometry}
\geometry{verbose,tmargin=1in,bmargin=1in,lmargin=1in,rmargin=1in}

\makeatletter
%%%%%%%%%%%%%%%%%%%%%%%%%%%%%% Textclass specific LaTeX commands.
\numberwithin{equation}{section}
\numberwithin{figure}{section}
\newlength{\lyxlabelwidth}      % auxiliary length 

%%%%%%%%%%%%%%%%%%%%%%%%%%%%%% User specified LaTeX commands.
\usepackage{fancyhdr}
\usepackage{lastpage}
\pagestyle{fancy}
\usepackage{enumitem}
\usepackage{ifthen}
\everymath=\expandafter{\the\everymath\displaystyle}
%\DeclareMathSizes{10}{10}{10}{10}
\usepackage{multicol}

\usepackage{tikz}
\usepackage{tikz-3dplot}
\usetikzlibrary{calc,arrows}

\usetikzlibrary{patterns}
\usetikzlibrary{plotmarks}
\usepackage{pgfplots}
\pgfplotsset{compat=newest}

\usepackage{calc}
\usepackage[nomessages]{fp}

\usepackage{datetime}
% Added by lyx2lyx
\setlength{\parskip}{\smallskipamount}
\setlength{\parindent}{0pt}

\makeatother

\usepackage{babel}
\begin{document}
\renewcommand{\author}{Dr.\,James Ashe}

\newcommand{\classNum}{336}

\newcommand{\classSec}{A}

\newcommand{\examType}{HW 1.5 Solutions}

\renewcommand{\date}{\formatdate{20}{9}{2013}}

%%First page's header%%%%%%%%%%%%%%%%%%%%%%
\fancypagestyle{firststyle} %%header settings for first pagr
{
	\fancyhead[L]{MTH \classNum{}(\classSec{}) \author{}\\
	\textbf{\examType{}} ~\date{}}
	\fancyhead[C]{}
	\fancyhead[R]{\textbf{Name:\hspace{4cm}}}
	\setlength{\headsep}{14pt}
}
%%%%%%%%%%%%%%%%%%%%%%%%%%%%%%%%%%%%%%%%%%

%%Next page's header%%%%%%%%%%%%%%%%%%%%%%%
\setlist{nolistsep}
\lhead{\classNum{}(\classSec{})} 
\chead{\textbf{Name:}} 
\cfoot{\thepage{}~of~\pageref{LastPage}} 
\setlength{\headsep}{5pt} 
%%%%%%%%%%%%%%%%%%%%%%%%%%%%%%%%%%%%%%%%%%%

%%Various settings; see the preamble for other settings%%%%%
%\setenumerate[0]{leftmargin=12pt,itemindent=0pt,labelwidth=0pt}
\setlist{topsep=.5pt}
\renewcommand{\labelenumi}{\textbf{\arabic{enumi}.}}
\renewcommand{\labelenumii}{\textbf{\alph{enumii}.}}
\thispagestyle{firststyle}
%%%%%%%%%%%%%%%%%%%%%%%%%%%%%%%%%%%%%%%%%%

\newcommand{\xmax}{0}
\newcommand{\xmin}{-\xmax}
\newcommand{\xstep}{0}
\newcommand{\ymax}{0}
\newcommand{\ymin}{-\ymax}
\newcommand{\ystep}{0} 
\newcommand{\dspace}{\vspace{1cm}}
\newcommand{\xa}{0}
\newcommand{\xb}{0}
\newcommand{\ya}{1}
\newcommand{\yb}{1}
\begin{enumerate}[resume]
\item [\textbf{3.}]

\begin{enumerate}[resume]
\item [\textbf{b.}]
\[
\begin{bmatrix}\begin{array}{rrr|r}
1 & 1 & 1 & b_{1}\\
-1 & 1 & 2 & b_{2}\\
1 & 3 & 4 & b_{3}
\end{array}\end{bmatrix}\rightsquigarrow\begin{bmatrix}\begin{array}{rrr|r}
1 & 1 & 1 & b_{1}\\
0 & 2 & 3 & b_{1}+b_{2}\\
0 & 2 & 3 & b_{3}-b_{1}
\end{array}\end{bmatrix}\rightsquigarrow\begin{bmatrix}\begin{array}{rrr|r}
1 & 1 & 1 & b_{1}\\
0 & 2 & 3 & b_{1}+b_{2}\\
0 & 0 & 0 & b_{3}+3b_{1}
\end{array}\end{bmatrix}
\]
So the constraint equation is $2b_{1}+b_{2}-b_{3}=0$.
\end{enumerate}
\item [\textbf{4.}]To be an element of $V$ we must solve $A\mathbf{x}=\mathbf{b}$.

\begin{enumerate}[resume]
\item [\textbf{a.}]
\[
\begin{bmatrix}\begin{array}{rr|r}
-1 & 2 & b_{1}\\
2 & -4 & b_{2}\\
1 & -2 & b_{3}
\end{array}\end{bmatrix}\rightsquigarrow\begin{bmatrix}\begin{array}{rr|r}
1 & -2 & -b_{1}\\
0 & 0 & b_{2}+2b_{1}\\
0 & 0 & b_{3}+b_{1}
\end{array}\end{bmatrix}
\]
So then $\mathbf{b}$ must satisfy the constraint equations $b_{2}+2b_{1}=0$
and $b_{3}+b_{1}=0$.
\item [\textbf{c.}]
\[
\begin{bmatrix}\begin{array}{rrr|r}
1 & 0 & 2 & b_{1}\\
0 & 1 & -1 & b_{2}\\
1 & 1 & 1 & b_{3}\\
1 & 2 & 0 & b_{4}
\end{array}\end{bmatrix}\rightsquigarrow\begin{bmatrix}\begin{array}{rrrr}
1 & 0 & 2 & b_{1}\\
0 & 1 & -1 & b_{2}\\
0 & 0 & 0 & b_{3}-b_{3}-b_{1}\\
0 & 0 & 0 & b_{4}-2b_{2}-b_{1}
\end{array}\end{bmatrix}
\]
So then $\mathbf{b}$ must satisfy the constraint equations $b_{1}+b_{2}-b_{3}=0$
and $b_{1}+2b_{2}-b_{4}=0$.
\end{enumerate}
\item [\textbf{5}]Rewrite the parametric equation as a matrix $\left[A\mathbf{x}|\mathbf{b}\right]$.
The Cartesian equation will then be the constraint equation. 

\begin{enumerate}[resume]
\item [\textbf{a.}]
\[
\begin{bmatrix}\begin{array}{rr|r}
1 & 2 & b_{1}\\
-2 & 0 & b_{2}\\
-2 & -1 & b_{3}
\end{array}\end{bmatrix}\rightsquigarrow\begin{bmatrix}\begin{array}{rr|r}
1 & 2 & b_{1}\\
0 & 4 & b_{2}+2b_{1}\\
0 & 0 & 2b_{1}-3b_{3}+4b_{3}
\end{array}\end{bmatrix}
\]
So the constraint equation give the Cartesian equation, $2x_{1}-3x_{3}+4x_{3}=0$.
\end{enumerate}
\item [\textbf{10.}]We can give a counterexample which proves that the
statement is false. Let 
\[
A=\begin{bmatrix}\begin{array}{rr}
0 & 0\\
0 & 0\\
0 & 0
\end{array}\end{bmatrix}
\]
 so that $A\mathbf{x}=\mathbf{0}$ certainly has the solution $\mathbf{x}=\mathbf{0}$,
but 
\[
A\mathbf{x}=\begin{bmatrix}\begin{array}{r}
1\\
0\\
0
\end{array}\end{bmatrix}
\]
 has no solution.
\item [\textbf{15.}]

\begin{enumerate}[resume]
\item [\textbf{a.}]$\mathbf{a_{1}+a_{2}+\cdots a}_{n}=\mathbf{0}$ implies
that $A\mathbf{x}=\mathbf{0}$ has the nontrivial solution $\mathbf{x}=(1,1,\dots,1)$
where $\mathbf{x}\in\mathbb{R}^{n}$. Thus by Proposition 5.5 $A$
is singular which by definition means that $\mbox{rank}(A)<n$. 
\item [\textbf{b.}]As above, if there is a nonzero vector $\mathbf{c}=(c_{1},c_{2},\dots,c_{n})$
such that $\sum_{i=1}^{n}c_{i}\mathbf{a}_{i}=\mathbf{0}$, then $A\mathbf{x}=\mathbf{0}$
has the nontrivial solution $\mathbf{x}=\mathbf{c}$ implying that
$A$ is singular and that $\mbox{rank}(A)<n$.
\end{enumerate}
\end{enumerate}

\end{document}
